\subsection{Einsatz einer Diode zur Leistungsreduzierung}
% Alt: Aufgabe 8

\Aufgabe{
    \label{Aufg8}
    In einer einfachen Schaltung wird eine Diode zur Leistungsreduktion eingesetzt. Der Widerstand $R_\mathrm{L} = 1\,\mathrm{k\Omega}$ stellt eine Heizlast dar.
    Die Netzspannung beträgt $230\,\mathrm{V}$ Effektivwert bei einer Frequenz von $50\,\mathrm{Hz}$.

    \begin{figure}[H]
        \centering
        \input{Tikz/src/FigDiodeLeistungsreduzierung}\alttext{Damit die Leistungsreduzierung veranschaulicht werden kann wurde eine elektrische Schaltung mit drei Bauteilen erstellt. Auf der linken Seite befindet sich deine  Wechselspannungsquelle \(u_\mathrm{E}\), deren Spannungspfeil nach unten zeigt. In Reihe zur Spannungsquelle sind ein Widerstand \(R_\mathrm{L}\) und eine Diode \(D\) geschaltet. Neben der Diode befindet sich zusätzlich ein Spannungspfeil in Stromrichtung mit der Beschriftung \(U_\mathrm{D}\).}
        \label{fig:FigDiodeLeistungsreduzierung}
    \end{figure}

    \begin{itemize}
        \item[a)] Welche mittlere Leistung wird in $R_\mathrm{L}$ und in der Diode umgesetzt, wenn die Diode als idealer Gleichtrichter wirkt (d.h. nur eine Halbwelle durchlässt)?
        \item[b)] Wie groß wäre die mittlere Leistung bei überbrückter Diode (volle Netzspannung an $R_\mathrm{L}$)?
    \end{itemize}
}


\Loesung{
    \begin{itemize}
        \item[a)]
              Bei einem idealen Ventil: Nach dem Erreichen der Flussspannung verhält sich die Diode wie ein perfekter Leiter ohne Widerstand. \\
              Da nur eine Halbwelle der Wechselspannung durchgelassen wird, halbiert sich die Leistung im Lastwiderstand. Die Diode selbst setzt keine Leistung um:
              \begin{align*}
                  P_\mathrm{D} & = 0\,\mathrm{W}
              \end{align*}
              Die mittlere Leistung im Lastwiderstand beträgt:
              \begin{align*}
                  P_\mathrm{RL} & = \frac{1}{2} \cdot \frac{(U_\mathrm{eff})^2}{R_\mathrm{L}} = \frac{1}{2} \cdot \frac{(230\,\mathrm{V})^2}{1\,\mathrm{k\Omega}} = 26,45\,\mathrm{W}
              \end{align*}

        \item[b)]
              Bei überbrückter Diode liegt die volle Wechselspannung am Lastwiderstand an:
              \begin{align*}
                  P_\mathrm{RL} & = \frac{U_\mathrm{eff}^2}{R_\mathrm{L}} = \frac{230\,\mathrm{V}^2}{1\,\mathrm{k\Omega}} = 52,9\,\mathrm{W}
              \end{align*}
    \end{itemize}

    Siehe: Abschnitt \ref{sec:Einweggleichrichter}
}